\documentclass[11pt]{article}

\usepackage[margin=1in]{geometry}
\usepackage{amsmath, amssymb, mathtools, bm}
\usepackage{booktabs, array, longtable}
\usepackage{hyperref}

\newcommand{\R}{\mathbb{R}}
\newcommand{\E}{\mathbb{E}}
\newcommand{\Var}{\mathrm{Var}}
\newcommand{\Cov}{\mathrm{Cov}}
\newcommand{\diag}{\mathrm{diag}}
\newcommand{\trans}{\mathsf{T}}

\title{How \texttt{predict-all-59.R} Extrapolates to Unknown Locations}
\author{}
\date{April 29, 2026}

\begin{document}

\maketitle

\section{Purpose}

This note translates the current implementation in the prediction script
\texttt{predict-all-59.R} and the helper \texttt{predictY\_cont\_lambda()}
from \texttt{update\_params.R} into matrix notation.
The goal is to explain, in mathematical terms, how a model that has already
been fitted on observed AQS sites is pushed forward to unobserved CMAQ grid
cells.

For setting 59, the fitted model is loaded from
\texttt{results/us-all-full-59.RData}. In the current checkout, the prediction
script uses:
\[
n_s = 1089 \quad \text{observed sites}, \qquad
n_p = 2632 \quad \text{prediction cells}, \qquad
n_t = 31 \quad \text{days},
\]
\[
n_{\text{reps}} = 5000 \quad \text{posterior draws}, \qquad
K = 6 \quad \text{spatial partitions (knots)}.
\]

Setting 59 is a symmetric \(t\) model with \texttt{skew = FALSE}, so the skew
term is inactive:
\[
\lambda = 0, \qquad z_{k,t} = 0 \quad \text{for all } k,t.
\]
Therefore the prediction step is driven by the regression mean, the spatial
correlation matrix, and the partition-specific precision parameters.

\section{Observed and Unknown Locations}

Let
\[
S_o =
\begin{bmatrix}
(s^o_1)^\trans \\
\vdots \\
(s^o_{n_s})^\trans
\end{bmatrix}
\in \R^{n_s \times 2}
\]
be the observed AQS locations after removing sites with more than \(50\%\)
missing responses, and let
\[
S_p =
\begin{bmatrix}
(s^p_1)^\trans \\
\vdots \\
(s^p_{n_p})^\trans
\end{bmatrix}
\in \R^{n_p \times 2}
\]
be the unobserved grid cells in the southeast subdomain selected by the code.

The script computes three distance matrices:
\[
D_{22} \in \R^{n_s \times n_s}, \qquad
(D_{22})_{ab} = \lVert s^o_a - s^o_b \rVert_2,
\]
\[
D_{11} \in \R^{n_p \times n_p}, \qquad
(D_{11})_{ab} = \lVert s^p_a - s^p_b \rVert_2,
\]
\[
D_{12} \in \R^{n_p \times n_s}, \qquad
(D_{12})_{ab} = \lVert s^p_a - s^o_b \rVert_2.
\]

These matrices are the only spatial inputs used by the prediction step.

\section{CMAQ Covariates and Design Matrices}

The raw CMAQ field is centered and scaled once:
\[
\widetilde{C}(s,t)
=
\frac{C(s,t) - \bar{C}}{s_C},
\]
where \(\bar{C}\) and \(s_C\) are the global mean and standard deviation used
by the script.

For each day \(t \in \{1,\dots,n_t\}\), define the observed-site and
prediction-site design matrices
\[
X_{o,t}
=
\begin{bmatrix}
1 & \widetilde{C}(s^o_1,t) \\
\vdots & \vdots \\
1 & \widetilde{C}(s^o_{n_s},t)
\end{bmatrix}
\in \R^{n_s \times 2},
\]
\[
X_{p,t}
=
\begin{bmatrix}
1 & \widetilde{C}(s^p_1,t) \\
\vdots & \vdots \\
1 & \widetilde{C}(s^p_{n_p},t)
\end{bmatrix}
\in \R^{n_p \times 2}.
\]

The regression coefficient vector from posterior draw \(i\) is
\[
\beta_i =
\begin{bmatrix}
\beta_{0,i} \\
\beta_{1,i}
\end{bmatrix}
\in \R^2.
\]
Hence the code-level regression means are
\[
\mu^{o}_{i,t} = X_{o,t} \beta_i \in \R^{n_s},
\qquad
\mu^{p}_{i,t} = X_{p,t} \beta_i \in \R^{n_p}.
\]

\section{\texttt{simple.cov.sp()} in Mathematical Form}

The function \texttt{simple.cov.sp()} receives a distance matrix
\(D = (d_{ab})\) and returns a covariance matrix \(V = (V_{ab})\). Write
\[
\texttt{sp.par} = (\sigma^2, \rho),
\]
where \(\sigma^2\) is a variance scale and \(\rho\) is a range parameter.

The returned covariance is
\[
V_{ab} = C_0(d_{ab}) + \mathbf{1}(d_{ab} = 0)\bigl(\tau^2_{\text{fine}} +
\tau^2_{\text{err}}\bigr),
\]
where \(\tau^2_{\text{fine}} = \texttt{finescale.var}\),
\(\tau^2_{\text{err}} = \texttt{error.var}\), and \(C_0\) depends on
\texttt{sp.type}:
\[
C_0(d) =
\begin{cases}
\sigma^2 \exp(-d/\rho), & \text{exponential}, \\[4pt]
\sigma^2 \exp\!\bigl(-(d/\rho)^2\bigr), & \text{gaussian}, \\[4pt]
\sigma^2 \dfrac{2^{1-\nu}}{\Gamma(\nu)}
\left(\dfrac{d}{\rho}\right)^\nu
K_\nu\!\left(\dfrac{d}{\rho}\right), & \text{matern}, \\[12pt]
\sigma^2 \dfrac{2^{1-\nu}}{\Gamma(\nu)}
\left(\dfrac{2\sqrt{\nu}\, d}{\rho}\right)^\nu
K_\nu\!\left(\dfrac{2\sqrt{\nu}\, d}{\rho}\right), & \text{matern2}, \\[12pt]
\sigma^2 \left(1 - \dfrac{3}{2}\dfrac{d}{\rho}
+ \dfrac{1}{2}\left(\dfrac{d}{\rho}\right)^3\right)
\mathbf{1}(d < \rho), & \text{spherical}.
\end{cases}
\]

Here \(\nu > 0\) is the Mat\'ern smoothness parameter and \(K_\nu(\cdot)\) is
the modified Bessel function of the second kind.

For setting 59, the call is
\[
\texttt{simple.cov.sp}(D, \texttt{"matern"}, (\sigma^2,\rho_i), 0, \nu_i, 0)
\quad \text{with } \sigma^2 = 1.
\]
Therefore the script uses the unit-variance Mat\'ern correlation kernel
\[
R_i(d)
=
\frac{2^{1-\nu_i}}{\Gamma(\nu_i)}
\left(\frac{d}{\rho_i}\right)^{\nu_i}
K_{\nu_i}\!\left(\frac{d}{\rho_i}\right),
\qquad d > 0,
\]
with the diagonal convention \(R_i(0)=1\). The implementation sets the
diagonal through the limiting value at \(d=0\), which is why the code replaces
\texttt{NaN} values by \(\sigma^2 = 1\).

\section{From the Mat\'ern Kernel to the Draw-Specific Correlation Matrix}

For each posterior draw \(i \in \{1,\dots,5000\}\), the script extracts three
draw-specific spatial parameters:
\[
\rho_i, \qquad \nu_i, \qquad \gamma_i.
\]

Define the Mat\'ern correlation matrices
\[
R^{oo}_i = R_i(D_{22}) \in \R^{n_s \times n_s},
\qquad
R^{pp}_i = R_i(D_{11}) \in \R^{n_p \times n_p},
\qquad
R^{po}_i = R_i(D_{12}) \in \R^{n_p \times n_s}.
\]

The code then constructs the observed-site correlation matrix as
\[
C_i = (1-\gamma_i) I_{n_s} + \gamma_i R^{oo}_i.
\]
This is exactly what the code implements by first computing
\(\gamma_i R^{oo}_i\) and then forcing the diagonal to be \(1\).

Likewise, the prediction-side covariance pieces are
\[
S_{11,i} = (1-\gamma_i) I_{n_p} + \gamma_i R^{pp}_i,
\qquad
S_{12,i} = \gamma_i R^{po}_i.
\]
The matrix named \texttt{prec} in the code is
\[
\texttt{prec}_i = C_i^{-1}.
\]

\subsection*{Why \(\texttt{prec}_i\) changes with \(i\)}

The prediction loop runs over \(i = 1,\dots,n_{\text{reps}}\). In each draw,
\(\rho_i\), \(\nu_i\), and \(\gamma_i\) change, so the matrices
\[
R^{oo}_i, \quad R^{pp}_i, \quad R^{po}_i, \quad C_i, \quad C_i^{-1}
\]
all change as well.

In the current fitted object for setting 59, the posterior draws vary roughly
over the following ranges:
\[
\rho_i \text{ quantiles } (5\%, 50\%, 95\%) \approx
(0.310,\ 0.388,\ 0.490),
\]
\[
\nu_i \text{ quantiles } (5\%, 50\%, 95\%) \approx
(0.258,\ 0.288,\ 0.331),
\]
\[
\gamma_i \text{ quantiles } (5\%, 50\%, 95\%) \approx
(0.899,\ 0.924,\ 0.943).
\]

This matters because:
\begin{itemize}
\item Larger \(\rho_i\) makes spatial dependence decay more slowly with
distance, so distant observed sites can influence a prediction more strongly.
\item Larger \(\nu_i\) changes the local smoothness of the Mat\'ern field and
therefore changes the shape of the correlation matrix.
\item Larger \(\gamma_i\) pushes \(C_i\) farther away from the identity matrix,
which means the observed residuals are treated as more spatially coherent and
the kriging correction becomes stronger.
\end{itemize}

An important code-level detail is that \(\texttt{prec}_i = C_i^{-1}\) depends
only on \((\rho_i,\nu_i,\gamma_i)\), not on \(t\). Therefore, within a fixed
posterior draw \(i\), the same inverse correlation matrix is reused for all 31
days. Day-to-day differences then come from the CMAQ covariate values, the
partition assignments, the partition precisions, and the imputed observed
responses.

\section{Partition Membership and Precision Scaling}

Setting 59 uses \(K=6\) partitions. For each posterior draw \(i\) and day \(t\),
the fitted object contains a \(6 \times 2\) knot matrix
\[
\mathcal{K}_{i,t} \in \R^{6 \times 2}.
\]

The helper \texttt{mem(s, knots)} assigns each location to the nearest knot.
Hence the observed-site and prediction-site memberships are
\[
g^o_{i,t}(a) \in \{1,\dots,6\}, \qquad a=1,\dots,n_s,
\]
\[
g^p_{i,t}(b) \in \{1,\dots,6\}, \qquad b=1,\dots,n_p.
\]

The fitted object also contains partition-specific precisions
\[
\tau_{i,k,t} > 0,
\qquad k=1,\dots,6.
\]
These are expanded to site-level quantities by lookup:
\[
\tau^o_{i,t}(a) = \tau_{i,\, g^o_{i,t}(a),\, t},
\qquad
\tau^p_{i,t}(b) = \tau_{i,\, g^p_{i,t}(b),\, t}.
\]

Now define the diagonal scaling matrices
\[
D^o_{i,t} = \diag\!\bigl(\sqrt{\tau^o_{i,t}(1)}, \dots,
\sqrt{\tau^o_{i,t}(n_s)}\bigr),
\]
\[
\Sigma^p_{i,t} = \diag\!\left(
\frac{1}{\sqrt{\tau^p_{i,t}(1)}}, \dots,
\frac{1}{\sqrt{\tau^p_{i,t}(n_p)}}
\right).
\]
These are exactly the code objects
\texttt{taug.t = sqrt(taug[, t])} and
\texttt{siggp = 1 / sqrt(tau[gp, t])}.

\section{The Conditional Prediction Formula}

For posterior draw \(i\) and day \(t\), the script loads
\[
y^o_{i,t} \in \R^{n_s},
\]
where \(y^o_{i,t}\) is the imputed observed response vector stored in
\texttt{fit\$y}. This is not the raw observation vector; it is the completed
response vector carried forward from the fitting stage.

Because setting 59 is non-skew, the mean term is just
\[
\mu^o_{i,t} = X_{o,t}\beta_i,
\qquad
\mu^p_{i,t} = X_{p,t}\beta_i.
\]
The observed residual vector is
\[
r^o_{i,t} = y^o_{i,t} - \mu^o_{i,t}.
\]
The code then standardizes this residual by the partition precisions:
\[
u^o_{i,t} = D^o_{i,t} r^o_{i,t}.
\]

The central conditional covariance object is
\[
Q_i = S_{11,i} - S_{12,i} C_i^{-1} S_{12,i}^\trans,
\]
which is the Schur complement appearing in Gaussian conditional
distributions. In the code this is the object named \texttt{corp}.

The predictive mean at unknown locations is
\[
m^p_{i,t}
=
\mu^p_{i,t}
+
\Sigma^p_{i,t} S_{12,i} C_i^{-1} u^o_{i,t}.
\]
This is the line
\[
\texttt{mup <- xp.beta + siggp * s.12.22.inv \%*\% (taug.t * res[, t])}.
\]

The predictive covariance is
\[
V^p_{i,t}
=
\Sigma^p_{i,t} Q_i \Sigma^p_{i,t}.
\]

Therefore the actual prediction law used by the script is
\[
Y^p_{i,t} \mid y^o_{i,t}, \theta_i
\sim
\mathcal{N}\!\left(m^p_{i,t},\, V^p_{i,t}\right),
\]
where \(\theta_i\) denotes all parameters loaded from posterior draw \(i\).

\section{How the Script Draws One Realization}

Let \(L_i\) be a Cholesky factor such that
\[
Q_i = L_i L_i^\trans.
\]
The code computes such a factor from \texttt{corp}; if Cholesky fails, it falls
back to an eigen-based square root.

If \(\varepsilon_{i,t} \sim \mathcal{N}(0, I_{n_p})\), then the script draws
\[
Y^p_{i,t}
=
m^p_{i,t}
+
\Sigma^p_{i,t} L_i \varepsilon_{i,t}.
\]
This is the exact meaning of
\[
\texttt{yp[, t] <- mup + siggp * t.corp.sd.mtx \%*\% rnorm(np)}.
\]

So the script is not merely computing a kriging mean. It is drawing a full
posterior predictive sample for every day and every prediction location.

\section{What Changes Across the Posterior Draw Index}

Each posterior draw \(i\) changes the extrapolation in three coupled ways:
\begin{enumerate}
\item The regression term changes through \(\beta_i\), which alters the CMAQ
baseline \(\mu^p_{i,t}\).
\item The spatial borrowing pattern changes through
\((\rho_i,\nu_i,\gamma_i)\), which changes \(C_i\), \(S_{11,i}\), \(S_{12,i}\),
and therefore \(\texttt{prec}_i = C_i^{-1}\).
\item The scale of the response changes through the day-specific partition
precisions \(\tau_{i,k,t}\) and the day-specific knot assignments
\(\mathcal{K}_{i,t}\).
\end{enumerate}

Consequently, the predictive distribution at a fixed unobserved location is not
based on a single deterministic set of kriging weights. It is averaged over
5000 distinct spatial systems:
\[
\left\{
\beta_i,\rho_i,\nu_i,\gamma_i,\tau_{i,\cdot,\cdot},
\mathcal{K}_{i,\cdot},y^o_{i,\cdot}
\right\}_{i=1}^{5000}.
\]

\section{Final Stored Object}

For each posterior draw \(i\), each prediction cell \(b\), and each day \(t\),
the script stores one predictive sample:
\[
\texttt{y.pred}[i,b,t] = Y^p_{i,t}(b).
\]
Hence the saved object is
\[
\texttt{y.pred} \in \R^{5000 \times 2632 \times 31}.
\]

This means the output is a Monte Carlo representation of the full posterior
predictive distribution on the southeast grid, not a single map or a single
plug-in predictor.

\section{One-Sentence Summary}

The prediction code takes each posterior draw from the fitted model, rebuilds
the draw-specific spatial correlation matrix \(C_i\), inverts it to obtain
\(\texttt{prec}_i\), uses observed-site residuals to form a conditional
Gaussian update, rescales that update by partition-specific precisions, and
then samples ozone values on previously unobserved CMAQ grid cells.

\end{document}
